% GENERATED FILE. Edit canonical lessons, then run npm run generate:books.
% canonical-source-hash: fnv1a64:cc0088f1fdbdc453
% canonical-lessons: ML-C407-teliyikkuka, ML-C407-samsayam, ML-C407-anumanikkuka, ML-C407-yathartthyam, ML-C407-artthamakkuka

\chapter{Proof and Doubt}
\label{ch:ml-proof-and-doubt}
\hlchaptermodality{407}

\begin{chapteropening}
\textbf{By the end of this chapter:} \emph{I can say to prove, a doubt, to assume, reality and to imply.}

Complete the last lesson of chapter 407: I can say to prove, a doubt, to assume, reality and to imply.
\end{chapteropening}

\section[\texorpdfstring{\ml{തെളിയിക്കുക} (teḷiyikkuka)}{തെളിയിക്കുക (teḷiyikkuka)}]{\ml{തെളിയിക്കുക} (teḷiyikkuka) — to prove}

\label{lesson:ML-C407-teliyikkuka}



\begin{quote}
Before the new one: say the Malayalam for to convince, then the Malayalam for an outlook.
\end{quote}

\subsection*{You'll want to know: \ml{തെളിയിക്കുക}}
\textbf{\ml{തെളിയിക്കുക}} — \emph{teḷiyikkuka} — \textquotedblleft{}to prove\textquotedblright{}.

\begin{quote}
\ml{തന്റെ} \ml{വാദം} \ml{ശരിയാണെന്ന്} \ml{അവൻ} \ml{തെളിയിച്ചു}.
\end{quote}

\begin{quote}
\emph{tanṟe vādaṁ śariyāṇennŭ avan teḷiyiccu.}
\end{quote}

\begin{quote}
He proved that his argument was right.
\end{quote}

Past: \ml{തെളിയിച്ചു}, \emph{teḷiyiccu}.

The noun is \ml{തെളിവ്}, proof. Both belong to \ml{തെളിയുക}, to become clear.

\begin{cousinweb}
\textbf{\ml{തെളിയിക്കുക}} (\emph{teḷiyikkuka}) is native Dravidian, the causative of \ml{തെളിയുക}, \textquotedblleft{}to become clear\textquotedblright{} (DEDR 3433). DEDR lists Malayalam \emph{teḷivu}, \textquotedblleft{}clearness, proof\textquotedblright{}, in the same article: to prove is to make clear.
\end{cousinweb}

\begin{grammarlens}[title={What you've built}]
One more word for giving reasons and opinions.
\end{grammarlens}

\subsection*{Your turn}
\begin{itemize}
\raggedright
  \item \emph{Say it:} \emph{teḷiyikkuka}
  \item \emph{Say it:} \emph{teḷiyikkuka}, once more
  \item \emph{Say it:} \emph{tanṟe vādaṁ śariyāṇennŭ avan teḷiyiccu.} — He proved that his argument was right.
  \item \emph{Recall:} say the Malayalam for to convince, then the Malayalam for an outlook, then say \emph{teḷiyikkuka} again
\end{itemize}

\begin{tcolorbox}[breakable,colback=teal!4,colframe=teal!35!black,title={Before you move on}]
What does \ml{തെളിയിക്കുക} mean? (\textquotedblleft{}To prove\textquotedblright{}.) Say it once more.
\end{tcolorbox}
\section[\texorpdfstring{\ml{സംശയം} (saṁśayaṁ)}{സംശയം (saṁśayaṁ)}]{\ml{സംശയം} (saṁśayaṁ) — a query; a doubt}

\label{lesson:ML-C407-samsayam}



\begin{quote}
Before the new one: say the Malayalam for an outlook, then the Malayalam for to prove.
\end{quote}

\subsection*{You'll want to know: \ml{സംശയം}}
\textbf{\ml{സംശയം}} — \emph{saṁśayaṁ} — \textquotedblleft{}a query; a doubt\textquotedblright{}.

\begin{quote}
\ml{ക്ലാസിന്} \ml{ശേഷം} \ml{അവൾ} \ml{അധ്യാപകനോട്} \ml{ഒരു} \ml{സംശയം} \ml{ചോദിച്ചു}.
\end{quote}

\begin{quote}
\emph{klāsinŭ śēṣaṁ avaḷ adhyāpakanōṭŭ oru saṁśayaṁ cōdiccu.}
\end{quote}

\begin{quote}
After class she asked the teacher about a doubt.
\end{quote}

The verb is \ml{സംശയിക്കുക}, to suspect. Without a doubt: \ml{സംശയമില്ലാതെ}.

\begin{cousinweb}
\textbf{\ml{സംശയം}} (\emph{saṁśayaṁ}) is borrowed from Sanskrit \emph{saṁśaya}, \textquotedblleft{}danger\textquotedblright{}, and then \textquotedblleft{}doubt\textquotedblright{} (CDIAL 13030).
\end{cousinweb}

\begin{grammarlens}[title={What you've built}]
One more word for giving reasons and opinions.
\end{grammarlens}

\subsection*{Your turn}
\begin{itemize}
\raggedright
  \item \emph{Say it:} \emph{saṁśayaṁ}
  \item \emph{Say it:} \emph{saṁśayaṁ}, once more
  \item \emph{Say it:} \emph{klāsinŭ śēṣaṁ avaḷ adhyāpakanōṭŭ oru saṁśayaṁ cōdiccu.} — After class she asked the teacher about a doubt.
  \item \emph{Recall:} say the Malayalam for an outlook, then the Malayalam for to prove, then say \emph{saṁśayaṁ} again
\end{itemize}

\begin{tcolorbox}[breakable,colback=teal!4,colframe=teal!35!black,title={Before you move on}]
What does \ml{സംശയം} mean? (\textquotedblleft{}A query; a doubt\textquotedblright{}.) Say it once more.
\end{tcolorbox}
\section[\texorpdfstring{\ml{അനുമാനിക്കുക} (anumānikkuka)}{അനുമാനിക്കുക (anumānikkuka)}]{\ml{അനുമാനിക്കുക} (anumānikkuka) — to assume; to infer}

\label{lesson:ML-C407-anumanikkuka}



\begin{quote}
Before the new one: say the Malayalam for to prove, then the Malayalam for a query.
\end{quote}

\subsection*{You'll want to know: \ml{അനുമാനിക്കുക}}
\textbf{\ml{അനുമാനിക്കുക}} — \emph{anumānikkuka} — \textquotedblleft{}to assume; to infer\textquotedblright{}.

\begin{quote}
\ml{കാർ} \ml{ഇല്ലാത്തതിനാൽ} \ml{അവർ} \ml{പോയി} \ml{എന്ന്} \ml{ഞാൻ} \ml{അനുമാനിച്ചു}.
\end{quote}

\begin{quote}
\emph{kār illāttatināl avar pōyi ennŭ ñān anumāniccu.}
\end{quote}

\begin{quote}
Because the car was not there, I assumed that they had gone.
\end{quote}

Past: \ml{അനുമാനിച്ചു}, \emph{anumāniccu}.

What is assumed comes before \ml{എന്ന്}: … \ml{എന്ന്} \ml{അനുമാനിക്കുക}. The noun is \ml{അനുമാനം}, an inference.

\begin{cousinweb}
\textbf{\ml{അനുമാനിക്കുക}} (\emph{anumānikkuka}) is built from Sanskrit \emph{anumāna}, \textquotedblleft{}inference, thought\textquotedblright{} (CDIAL 328), from \emph{anu-}, \textquotedblleft{}after\textquotedblright{}, and \emph{mā}, \textquotedblleft{}to measure\textquotedblright{}, with the Malayalam verb ending -\ml{ിക്കുക}: to measure one thing by another.
\end{cousinweb}

\begin{grammarlens}[title={What you've built}]
One more word for giving reasons and opinions.
\end{grammarlens}

\subsection*{Your turn}
\begin{itemize}
\raggedright
  \item \emph{Say it:} \emph{anumānikkuka}
  \item \emph{Say it:} \emph{anumānikkuka}, once more
  \item \emph{Say it:} \emph{kār illāttatināl avar pōyi ennŭ ñān anumāniccu.} — Because the car was not there, I assumed that they had gone.
  \item \emph{Recall:} say the Malayalam for to prove, then the Malayalam for a query, then say \emph{anumānikkuka} again
\end{itemize}

\begin{tcolorbox}[breakable,colback=teal!4,colframe=teal!35!black,title={Before you move on}]
What does \ml{അനുമാനിക്കുക} mean? (\textquotedblleft{}To assume; to infer\textquotedblright{}.) Say it once more.
\end{tcolorbox}
\section[\texorpdfstring{\ml{യാഥാർത്ഥ്യം} (yāthārtthyaṁ)}{യാഥാർത്ഥ്യം (yāthārtthyaṁ)}]{\ml{യാഥാർത്ഥ്യം} (yāthārtthyaṁ) — reality}

\label{lesson:ML-C407-yathartthyam}



\begin{quote}
Before the new one: say the Malayalam for a query, then the Malayalam for to assume.
\end{quote}

\subsection*{You'll want to know: \ml{യാഥാർത്ഥ്യം}}
\textbf{\ml{യാഥാർത്ഥ്യം}} — \emph{yāthārtthyaṁ} — \textquotedblleft{}reality\textquotedblright{}.

\begin{quote}
\ml{സ്വപ്നവും} \ml{യാഥാർത്ഥ്യവും} \ml{വ്യത്യസ്തമാണ്}.
\end{quote}

\begin{quote}
\emph{svapnavuṁ yāthārtthyavuṁ vyatyastamāṇŭ.}
\end{quote}

\begin{quote}
Dreams and reality are different.
\end{quote}

A learned word; in talk people also say \ml{സത്യം}, the truth.

\begin{cousinweb}
\textbf{\ml{യാഥാർത്ഥ്യം}} (\emph{yāthārtthyaṁ}) is borrowed from Sanskrit \emph{yāthārthya}, \textquotedblleft{}agreement with the truth, the real meaning\textquotedblright{} (Monier-Williams), made from \emph{yathārtha}, \textquotedblleft{}as the thing is\textquotedblright{}.
\end{cousinweb}

\begin{grammarlens}[title={What you've built}]
One more word for giving reasons and opinions.
\end{grammarlens}

\subsection*{Your turn}
\begin{itemize}
\raggedright
  \item \emph{Say it:} \emph{yāthārtthyaṁ}
  \item \emph{Say it:} \emph{yāthārtthyaṁ}, once more
  \item \emph{Say it:} \emph{svapnavuṁ yāthārtthyavuṁ vyatyastamāṇŭ.} — Dreams and reality are different.
  \item \emph{Recall:} say the Malayalam for a query, then the Malayalam for to assume, then say \emph{yāthārtthyaṁ} again
\end{itemize}

\begin{tcolorbox}[breakable,colback=teal!4,colframe=teal!35!black,title={Before you move on}]
What does \ml{യാഥാർത്ഥ്യം} mean? (\textquotedblleft{}Reality\textquotedblright{}.) Say it once more.
\end{tcolorbox}
\section[\texorpdfstring{\ml{അർത്ഥമാക്കുക} (artthamākkuka)}{അർത്ഥമാക്കുക (artthamākkuka)}]{\ml{അർത്ഥമാക്കുക} (artthamākkuka) — to imply; to mean}

\label{lesson:ML-C407-artthamakkuka}



\begin{quote}
Before the new one: say the Malayalam for to assume, then the Malayalam for reality.
\end{quote}

\subsection*{You'll want to know: \ml{അർത്ഥമാക്കുക}}
\textbf{\ml{അർത്ഥമാക്കുക}} — \emph{artthamākkuka} — \textquotedblleft{}to imply; to mean\textquotedblright{}.

\begin{quote}
\ml{നീ} \ml{എന്താണ്} \ml{അർത്ഥമാക്കുന്നത്}?
\end{quote}

\begin{quote}
\emph{nī entāṇŭ artthamākkunnatŭ?}
\end{quote}

\begin{quote}
What do you mean?
\end{quote}

Past: \ml{അർത്ഥമാക്കി}, \emph{artthamākki}.

Built on \ml{അർത്ഥം}, a meaning, with \ml{ആക്കുക}, to make.

\begin{cousinweb}
\textbf{\ml{അർത്ഥമാക്കുക}} (\emph{artthamākkuka}) is a Malayalam compound: \ml{അർത്ഥം}, \textquotedblleft{}a meaning\textquotedblright{}, from Sanskrit \emph{artha}, \textquotedblleft{}aim, cause, thing\textquotedblright{} (CDIAL 638), and \ml{ആക്കുക}, \textquotedblleft{}to make\textquotedblright{}, native Dravidian (DEDR 333).
\end{cousinweb}

\begin{grammarlens}[title={What you've built}]
One more word for giving reasons and opinions.
\end{grammarlens}

\subsection*{Your turn}
\begin{itemize}
\raggedright
  \item \emph{Say it:} \emph{artthamākkuka}
  \item \emph{Say it:} \emph{artthamākkuka}, once more
  \item \emph{Say it:} \emph{nī entāṇŭ artthamākkunnatŭ?} — What do you mean?
  \item \emph{Recall:} say the Malayalam for to assume, then the Malayalam for reality, then say \emph{artthamākkuka} again
\end{itemize}

\begin{tcolorbox}[breakable,colback=teal!4,colframe=teal!35!black,title={Before you move on}]
What does \ml{അർത്ഥമാക്കുക} mean? (\textquotedblleft{}To imply; to mean\textquotedblright{}.) Say it once more.
\end{tcolorbox}
